% =====================================================================
%  Wehrkalender -- Vorschlag für die Exportschnittstelle von MP Feuer
%  Vorlage: deb-support (debsupport.sty, siehe latex/README.md)
%  Übersetzen mit:  pdflatex schnittstelle_mp_feuer.tex   (2x, wegen Seitenzahl)
% =====================================================================
\documentclass[a4paper,fontsize=10pt]{scrartcl}
\usepackage[bericht]{debsupport}

\usepackage{listings}
\newcolumntype{L}[1]{>{\raggedright\arraybackslash}p{#1}}
\colorlet{grauhell}{dshell}
\colorlet{grau}{dsgrau}

\lstset{
  basicstyle=\ttfamily\small,
  backgroundcolor=\color{grauhell},
  frame=single, rulecolor=\color{grau!40},
  framesep=4pt, xleftmargin=6pt, xrightmargin=6pt,
  columns=fullflexible, keepspaces=true,
  breaklines=true, showstringspaces=false,
  literate={ä}{{\"a}}1 {ö}{{\"o}}1 {ü}{{\"u}}1 {Ä}{{\"A}}1 {Ö}{{\"O}}1 {Ü}{{\"U}}1 {ß}{{\ss}}1 {–}{{--}}1
}
\newcommand{\code}[1]{\texttt{#1}}

\renewcommand*{\dsdokumentart}{Schnittstellenvorschlag}
\renewcommand*{\dskurztitel}{Export aus MP Feuer für den Wehrkalender}
\renewcommand*{\dsbeleg}{}
\setlength{\emergencystretch}{3em}

\begin{document}

\dstitel{Export aus MP Feuer}
        {Vorschlag für Inhalt und Format der Schnittstelle zum Wehrkalender}
        {Datum:   & \today\\
         Autor:   & \dsInhaber\\
         Projekt: & Wehrkalender (erster Einsatz: FF Kreuzau)\\
         Version: & Vorschlag 1}

\tableofcontents

% ------------------------------------------------------------------
\section{Zweck}

Der \emph{Wehrkalender} bringt die Termine einer Feuerwehr aus MP Feuer nur
lesend auf die Smartphones, Tablets und in das Outlook der Mitglieder. Jedes
Mitglied sieht die Kalender seiner Einheiten, Gruppen und Lehrgänge. Der
Kalenderserver übernimmt dazu regelmäßig Mitglieder, Zuordnungen und Termine
aus MP Feuer. Es werden keine Daten nach MP Feuer zurückgeschrieben.

Dieses Dokument beschreibt, welche Daten der Wehrkalender braucht und in
welcher Form er sie am einfachsten verarbeiten kann. Es ist ein Vorschlag:
Wenn MP Feuer Daten bereits in anderer Form bereitstellt, passen wir uns an,
solange die Inhalte aus Abschnitt~\ref{sec:muss} enthalten sind. Der
Quellcode des Wehrkalenders ist öffentlich
(\code{https://github.com/brami/Wehrkalender}, AGPL-3.0).

% ------------------------------------------------------------------
\section{Grundsätze}\label{sec:muss}

\begin{enumerate}
  \item \textbf{Vollständiger Stand statt Änderungen:} Jeder Abruf liefert den
        vollständigen aktuellen Stand aller Mitglieder und Termine des
        vereinbarten Zeitraums. Was fehlt, gilt als entfernt. Änderungslisten
        (\glqq geändert seit\grqq) sind nicht nötig.
  \item \textbf{Stabile Kennungen:} Mitglieder, Bereiche (Einheiten, Gruppen,
        Lehrgänge) und Termine haben je eine eindeutige Kennung, die sich nie
        ändert -- auch nicht bei einer Namensänderung, einem Umzug des Termins
        oder einer Umbenennung -- und die nie neu vergeben wird.
  \item \textbf{Gleiche Schlüssel:} Termine und Mitgliedschaften verwenden
        dieselben Schlüssel für Bereiche. Nur so lässt sich ein Termin dem
        Kalender zuordnen, den die Mitglieder dieses Bereichs sehen.
  \item \textbf{Ganz oder gar nicht:} Ein Abruf ist vollständig oder schlägt
        mit Fehlercode fehl. Ein halber Export würde beim Abgleich Mitglieder
        und Termine löschen. Zur Kontrolle enthält jeder Export die Anzahl der
        Einträge.
  \item \textbf{Datensparsamkeit:} Nur die unten genannten Felder. Keine
        Geburtsdaten, Anschriften, Telefonnummern, Qualifikationen oder
        Gesundheitsdaten.
\end{enumerate}

% ------------------------------------------------------------------
\section{Bereitstellung}

Bevorzugt ist ein \textbf{Abruf über HTTPS} durch den Kalenderserver:

\begin{center}
\begin{tabularx}{\textwidth}{L{4.2cm}X}
\toprule
\textbf{Punkt} & \textbf{Vorschlag} \\
\midrule
Adressen & \code{GET …/export/mitglieder}, \code{GET …/export/bereiche},
  \code{GET …/export/termine} (oder ein gemeinsamer Abruf mit allen drei
  Teilen) \\
Anmeldung & API-Schlüssel je Feuerwehr im Header
  (\code{Authorization: Bearer …}); nur Lesezugriff auf den Export \\
Format & JSON, UTF-8; Zeitangaben ISO~8601 \\
Häufigkeit & Abruf alle 15 Minuten; bei unverändertem Stand Antwort
  \code{304 Not Modified} über \code{ETag}/\code{If-None-Match} (wünschenswert,
  nicht zwingend) \\
Fehler & HTTP-Statuscode $\neq$ 200 und kein Inhalt; niemals ein
  unvollständiger Export mit 200 \\
Umfang & eine Feuerwehr je API-Schlüssel; bei rund 200 Mitgliedern und
  einigen Hundert Terminen ohne Seitenaufteilung \\
\bottomrule
\end{tabularx}
\end{center}

\textbf{Alternativ} legt MP Feuer die drei Dateien regelmäßig per SFTP auf dem
Kalenderserver ab (eigener Zugang nur per SSH-Schlüssel). Damit keine halb
geschriebene Datei gelesen wird: zuerst unter \code{*.part} hochladen und erst
danach umbenennen.

% ------------------------------------------------------------------
\section{Inhalt}

Alle drei Teile beginnen mit einem Kopf; die eigentlichen Einträge stehen in
\code{eintraege}:

\begin{lstlisting}
{
  "feuerwehr": "FF Kreuzau",
  "erstellt": "2027-01-15T06:00:00+01:00",
  "anzahl": 214,
  "eintraege": [ … ]
}
\end{lstlisting}

\subsection{Mitglieder}

Alle aktiven Mitglieder einschließlich Jugend- und Kinderfeuerwehr.
Ausgeschiedene Mitglieder fehlen im Export.

\begin{center}
\begin{tabularx}{\textwidth}{L{3.6cm}L{1.6cm}X}
\toprule
\textbf{Feld} & \textbf{Pflicht} & \textbf{Bedeutung} \\
\midrule
\code{id} & ja & unveränderliche Kennung des Mitglieds \\
\code{vorname}, \code{nachname} & ja & für Benutzername und Anrede \\
\code{email} & ja, falls vorhanden & für Zugang und Benachrichtigungen;
  leer, wenn keine hinterlegt ist. Mehrere Mitglieder dürfen dieselbe
  Adresse haben (z.\,B.\ Geschwister). \\
\code{mitgliedschaften} & ja & Liste der Schlüssel aller Bereiche, denen das
  Mitglied angehört (\code{schluessel}), optional mit \code{rolle}, z.\,B.\
  \code{teilnehmer}, \code{betreuer}, \code{ausbilder} \\
\bottomrule
\end{tabularx}
\end{center}

\begin{lstlisting}
{ "id": "4711",
  "vorname": "Jürgen", "nachname": "Müller",
  "email": "juergen.mueller@example.org",
  "mitgliedschaften": [
    { "schluessel": "LG-DROVE" },
    { "schluessel": "JF-GR1", "rolle": "betreuer" },
    { "schluessel": "LEHRGANG-KREIS-MASCH-2027-2",
      "rolle": "teilnehmer" } ] }
\end{lstlisting}

\subsection{Bereiche}

Alle Einheiten, Gruppen, Abteilungen und Lehrgänge -- auch solche ohne
Termine. Daraus legt der Kalenderserver die Kalender an.

\begin{center}
\begin{tabularx}{\textwidth}{L{3.6cm}L{1.6cm}X}
\toprule
\textbf{Feld} & \textbf{Pflicht} & \textbf{Bedeutung} \\
\midrule
\code{schluessel} & ja & unveränderlicher Schlüssel, z.\,B.\
  \code{LG-DROVE}, \code{JF-GR1} \\
\code{bezeichnung} & ja & Anzeigename, z.\,B.\ \glqq LG Drove\grqq \\
\code{art} & ja & \code{einheit}, \code{gruppe}, \code{abteilung},
  \code{lehrgang} oder \code{sonstiges} \\
\code{uebergeordnet} & nein & Schlüssel des übergeordneten Bereichs \\
\multicolumn{3}{@{}l}{\emph{nur bei \code{art = lehrgang}:}} \\
\code{ebene} & ja & \code{gemeinde}, \code{kreis}, \code{bezirk} oder
  \code{land} \\
\code{lehrgang} & ja & Bezeichnung \emph{genau wie beim Ausrichter},
  z.\,B.\ \glqq Maschinist 2/2027\grqq\ (mit laufender Nummer und Jahr) \\
\code{beginn}, \code{ende} & ja & erster und letzter Lehrgangstag
  (\code{JJJJ-MM-TT}) \\
\code{ausrichter\_kennung} & nein & Lehrgangsnummer des Ausrichters, falls
  vorhanden \\
\bottomrule
\end{tabularx}
\end{center}

\begin{lstlisting}
{ "schluessel": "LEHRGANG-KREIS-MASCH-2027-2",
  "bezeichnung": "Maschinist 2/2027",
  "art": "lehrgang", "ebene": "kreis",
  "lehrgang": "Maschinist 2/2027",
  "beginn": "2027-06-05", "ende": "2027-06-26" }
\end{lstlisting}

\textbf{Warum diese Lehrgangsangaben?} Die Termine der Kreislehrgänge
veröffentlicht der Kreis als Kalenderdatei. MP Feuer kennt die Teilnehmer.
Mit Ebene, Bezeichnung und Beginn ordnet der Wehrkalender einen Lehrgang aus
MP Feuer automatisch den Terminen des Kreises zu und zeigt sie nur den
Teilnehmern -- ohne dass die Termine in MP Feuer doppelt gepflegt werden
müssen. Fehlt die laufende Nummer in der Bezeichnung, genügt der Beginn zur
eindeutigen Zuordnung. Wo MP Feuer die Termine eines Lehrgangs selbst führt
(z.\,B.\ auf Gemeindeebene), kommen sie wie alle anderen Termine über den
Export.

\subsection{Termine}

Zeitraum: Termine ab zwölf Monaten in der Vergangenheit und alle künftigen.
Serientermine werden als einzelne Termine geliefert.

\begin{center}
\begin{tabularx}{\textwidth}{L{3.6cm}L{1.6cm}X}
\toprule
\textbf{Feld} & \textbf{Pflicht} & \textbf{Bedeutung} \\
\midrule
\code{id} & ja & unveränderliche Kennung des Termins; bei Serien je
  Einzeltermin \\
\code{bereich} & ja & Schlüssel des Bereichs (wie in den Mitgliedschaften) \\
\code{titel} & ja & \\
\code{beginn}, \code{ende} & ja & mit Uhrzeit und Zeitzone, z.\,B.\
  \code{2027-03-04T19:00:00+01:00}; bei ganztägigen Terminen nur das Datum,
  \code{ende} ist der letzte Tag \\
\code{ganztaegig} & ja & \code{true}/\code{false} \\
\code{ort}, \code{beschreibung} & nein & \\
\code{abgesagt} & ja & \code{true}, solange ein abgesagter Termin noch
  angezeigt werden soll; danach entfällt er aus dem Export \\
\code{geaendert} & nein & Zeitpunkt der letzten Änderung \\
\bottomrule
\end{tabularx}
\end{center}

\begin{lstlisting}
{ "id": "T-98765",
  "bereich": "LG-DROVE",
  "titel": "Übungsdienst",
  "beginn": "2027-03-04T19:00:00+01:00",
  "ende": "2027-03-04T21:00:00+01:00",
  "ganztaegig": false,
  "ort": "Gerätehaus Drove",
  "beschreibung": "Thema: Technische Hilfe",
  "abgesagt": false,
  "geaendert": "2027-02-20T10:12:00+01:00" }
\end{lstlisting}

% ------------------------------------------------------------------
\section{Was der Wehrkalender damit macht}

Zur Einordnung, nicht Teil der Schnittstelle:

\begin{itemize}
  \item Mitglieder mit gültiger E-Mail-Adresse erhalten einen Zugang;
        Mitglieder, die im Export fehlen, werden mit allen Daten gelöscht.
  \item Je Bereich gibt es einen Kalender; jedes Mitglied sieht genau die
        Kalender seiner Mitgliedschaften und kann sie einzeln abschalten.
  \item Bricht ein Export plötzlich stark ein (weniger als 80\,\% der
        Mitglieder oder Termine des letzten Abrufs), wird nichts übernommen
        und die Administratoren werden gewarnt.
  \item Kommt länger als 24 Stunden kein gültiger Export an, werden die
        Administratoren ebenfalls gewarnt.
\end{itemize}

% ------------------------------------------------------------------
\section{Fragen und nächste Schritte}

\begin{enumerate}
  \item Welche Teile dieses Vorschlags kann MP Feuer umsetzen, und wo weicht
        das vorhandene Datenmodell ab (z.\,B.\ Lehrgänge, Rollen,
        Zuordnungen von Terminen zu Bereichen)?
  \item Werden Lehrgänge in MP Feuer mit der Bezeichnung des Ausrichters
        angelegt, und sind Ebene sowie Beginn erfasst?
  \item Abruf über HTTPS oder Ablage per SFTP? Wie erfolgt die Anmeldung?
  \item Bereitstellung eines Testzugangs bzw.\ anonymisierter Beispieldaten
        der FF Kreuzau.
  \item Fallen Kosten an oder ist ein bestimmtes Lizenzmodell nötig?
\end{enumerate}

Für Rückfragen und eine gemeinsame Abstimmung stehen wir gern zur Verfügung.

\end{document}
